% =====================================================================
%  4.4  Atividade 3.2.2.4 -- Esgotamento descentralizado em comunidades
% =====================================================================
\subsection{Atividade 3.2.2.4 -- Instalar sistemas descentralizados de esgotamento sanitário em comunidades rurais}
\label{subsec:esgotamento-comunidades}

\subsubsection{Objetivo e fundamentação}

Instalar sistemas descentralizados de esgotamento sanitário em 70\% dos
domicílios das comunidades rurais atendidas pelos sistemas coletivos de
abastecimento de água (\atividade{3.2.2.3}), totalizando 4.067 domicílios.
A execução terá início após a seleção e a mobilização das prefeituras e
comunidades e após o diagnóstico da disposição final do esgoto de cada
domicílio, obtido no trabalho técnico-social (\atividade{3.2.2.2})
\cite{mop2026}.

\begin{fichaatividade}{Ficha-síntese da Atividade 3.2.2.4}{qua:ficha-3224}
  \campo{Código}{3.2.2.4}
  \campo{Tipo de atividade}{Serviço e obra}
  \campo{Forma de execução}{Contratação pública (licitação)}
  \campo{Fonte de recursos}{BIRD}
  \campo{Responsável}{IAT / IDR-Paraná}
  \campo{Dependências institucionais}{Prefeituras}
  \campo{Prazo estimado}{Ano 2 ao ano 5}
  \campo{Produto esperado}{4.067 sistemas implantados}
  \campo{Unidade de medida}{Sistema implantado}
  \campo{Evidência de comprovação}{Relatório técnico aprovado e comprovante de pagamento}
  \campo{Periodicidade}{Por lote}
  \campo{Validação}{IAT}
\end{fichaatividade}

\subsubsection{Procedimentos}
\label{subsubsec:proc-esgotamento}

Os procedimentos, representados na \autoref{fig:fluxograma-esgotamento}, são
organizados em cinco fases.

\paragraph{Fase 1 -- Contratação.}
\begin{enumerate}[label=\alph*)]
  \item \textbf{Termo de Referência e edital}: elaboração e aprovação dos TR
        para aquisição das soluções descentralizadas e para contratação das
        obras civis de instalação; abertura dos editais; realização das
        licitações; contratação das empresas.
\end{enumerate}

\paragraph{Fase 2 -- Projeto e instalação.}
\begin{enumerate}[label=\alph*), resume]
  \item \textbf{projeto individual por propriedade}: a empresa contratada
        elaborará o projeto de cada residência, considerando as
        características do terreno, o número de moradores e a solução de
        tratamento adequada;
  \item \textbf{instalação do sistema}: execução das obras civis de
        instalação das estruturas e do dispositivo de esgotamento sanitário
        domiciliar.
\end{enumerate}

\paragraph{Fase 3 -- Verificação técnica e capacitação.}
\begin{enumerate}[label=\alph*), resume]
  \item \textbf{atestado de funcionalidade e garantia}: verificação de que o
        sistema funciona conforme os requisitos definidos;
  \item \textbf{decisão técnica}: se o sistema atender às condições
        técnicas, segue para a etapa seguinte; caso contrário, serão feitas
        correções antes da validação final;
  \item \textbf{capacitação dos proprietários}: concomitantemente à
        instalação, os beneficiários serão orientados sobre o funcionamento
        do sistema, as práticas adequadas de uso e a manutenção preventiva,
        com registro e comprovação.
\end{enumerate}

\paragraph{Fase 4 -- Medição e comprovação.}
\begin{enumerate}[label=\alph*), resume]
  \item \textbf{relatórios de medição}: registro técnico do andamento e da
        execução das intervenções nas propriedades;
  \item \textbf{comprovação da implementação}: documentação da elaboração e
        da execução dos projetos, com evidência das intervenções;
  \item \textbf{assinatura do proprietário}: confirmação, pelo beneficiário,
        da execução das melhorias em sua propriedade.
\end{enumerate}

\paragraph{Fase 5 -- Verificação independente, pagamento e encerramento.}
\begin{enumerate}[label=\alph*), resume]
  \item \textbf{auditoria documental}: verificação da documentação gerada;
  \item \textbf{verificação \textit{in loco}}: visitas técnicas para
        confirmar a execução das intervenções;
  \item \textbf{relatório de medição final}: consolidação das informações
        técnicas da execução;
  \item \textbf{amostragem \textit{in loco}}: avaliação, por amostragem, da
        qualidade das intervenções;
  \item \textbf{ordenação das despesas}: validação financeira e pagamento
        dos serviços;
  \item \textbf{finalização}: encerramento formal do processo.
\end{enumerate}

\begin{figure}[htbp]
  \centering
  \caption{Fluxograma de instalação dos sistemas descentralizados de esgotamento sanitário (\atividade{3.2.2.4} e \atividade{3.2.2.5})}
  \label{fig:fluxograma-esgotamento}
  % Fluxograma de instalação dos sistemas descentralizados de esgotamento
% sanitário (Atividades 3.2.2.4 e 3.2.2.5)
\begin{tikzpicture}[node distance=0.45cm]

  % --- Coluna 1: contratação, projeto, instalação e verificação -----
  \node[etapa=4.2cm] (tr) {TR, edital e licitação};
  \node[etapa=4.2cm, below=of tr] (contr) {Contratação das empresas};
  \node[etapa=4.2cm, below=of contr] (proj) {Projeto individual da propriedade};
  \node[etapa=4.2cm, below=of proj] (inst) {Instalação do sistema\\(obras civis)};
  \node[etapa=4.2cm, below=of inst] (atest) {Atestar funcionalidade e garantia};
  \node[decisao, below=0.55cm of atest] (dec) {Atende aos\\requisitos?};
  \node[correcao, left=0.9cm of dec] (corr) {Correções técnicas};

  % --- Coluna 2: capacitação, medição, verificação e pagamento ------
  \node[etapa=4.6cm, right=3.2cm of tr] (cap) {Capacitação dos proprietários\\e registro};
  \node[etapa=4.6cm, below=of cap] (med) {Relatórios de medição};
  \node[etapa=4.6cm, below=of med] (comp) {Comprovação da implementação\\e assinatura do proprietário};
  \node[etapa=4.6cm, below=of comp] (aud) {Auditoria documental e\\verificação \textit{in loco}};
  \node[etapa=4.6cm, below=of aud] (medf) {Relatório de medição final\\e amostragem \textit{in loco}};
  \node[etapa=4.6cm, below=of medf] (ord) {Ordenação das despesas\\(pagamento)};
  \node[marco=4.6cm, below=of ord] (fim) {Finalização};

  % --- Ligações ----------------------------------------------------
  \draw[seta] (tr) -- (contr);
  \draw[seta] (contr) -- (proj);
  \draw[seta] (proj) -- (inst);
  \draw[seta] (inst) -- (atest);
  \draw[seta] (atest) -- (dec);
  \draw[seta] (dec) -- node[rotulo, above] {Não} (corr);
  \draw[seta] (corr) |- (atest.west);
  \draw[seta] (dec.east) -- node[rotulo, above] {Sim} ++(1.1,0) coordinate (c1)
              |- (cap.west);
  \draw[seta] (cap) -- (med);
  \draw[seta] (med) -- (comp);
  \draw[seta] (comp) -- (aud);
  \draw[seta] (aud) -- (medf);
  \draw[seta] (medf) -- (ord);
  \draw[seta] (ord) -- (fim);

  % --- Fases (rótulos laterais) -------------------------------------
  \begin{scope}[on background layer]
    \node[fill=cinzaclaro!60, rounded corners, inner sep=5pt, fit=(tr)(contr)] (f1) {};
    \node[fill=cinzaclaro!60, rounded corners, inner sep=5pt, fit=(proj)(inst)] (f2) {};
    \node[fill=cinzaclaro!60, rounded corners, inner sep=5pt, fit=(atest)(dec)(corr)] (f3) {};
    \node[fill=cinzaclaro!60, rounded corners, inner sep=5pt, fit=(cap)(med)(comp)] (f4) {};
    \node[fill=cinzaclaro!60, rounded corners, inner sep=5pt, fit=(aud)(medf)(ord)(fim)] (f5) {};
  \end{scope}
  \node[font=\scriptsize\bfseries, text=azulpsh, anchor=south, rotate=90] at (f1.west) {Fase 1};
  \node[font=\scriptsize\bfseries, text=azulpsh, anchor=south, rotate=90] at (f2.west) {Fase 2};
  \node[font=\scriptsize\bfseries, text=azulpsh, anchor=north, rotate=90] at (f4.east) {Fases 3--4};
  \node[font=\scriptsize\bfseries, text=azulpsh, anchor=north, rotate=90] at (f5.east) {Fase 5};
  \node[font=\scriptsize\bfseries, text=azulpsh, anchor=north west] at (f3.north west) {Fase 3};
\end{tikzpicture}

  \fonte{Elaborado pelo IAT (2026) com base em \citeonline{mop2026}.}
\end{figure}

\pendencia{O MOP não indica a tecnologia de tratamento prevista (por exemplo,
tanque séptico com filtro anaeróbio ou sumidouro) nem as normas técnicas de
projeto. Sugere-se citar a ABNT NBR 7229 (tanques sépticos) e a ABNT NBR 13969
(unidades de tratamento complementar e disposição final).}

\pendencia{Não há quadro de etapas com prazos, produtos intermediários e
evidências, como nas \atividade{3.2.2.1} e \atividade{3.2.2.3}. Também não
estão definidos o tamanho da amostra da verificação \textit{in loco} nem a
diferença entre ``verificação \textit{in loco}'' e ``amostragem \textit{in
loco}''.}

\subsubsection{Interfaces}

\begin{itemize}
  \item \atividade{3.2.2.2}: fornece o diagnóstico domiciliar que orienta os
        projetos individuais.
  \item \atividade{3.2.2.3}: define o universo de domicílios (70\% dos
        atendidos com abastecimento de água).
  \item \atividade{3.2.2.5}: adota o mesmo procedimento
        (\autoref{subsec:esgotamento-idr}).
\end{itemize}
